\section{Conclusion}\label{sec:conclusion}


%
The current deployment of feature attribution methods in high-stakes settings such as medicine and law
demands a rigorous understanding of their performance.
In this work, we characterized conditions under which common \fielddefnames{} provably are unreliable when used to infer \behaviourname{} for learned \modelnames{}. Our work concludes that using feature attribution as is currently prescribed in the literature need not improve on random guessing at inferring model behavior.
The implications of our work are two-fold.
First, it is necessary to build more structure into these methods in order to improve upon the current performance guarantees. We show that a simple brute-force method (easy to define but computationally expensive) is guaranteed to work for \textit{some \downtasknames{}}.
As a result, our second implication is that defining the \downtaskname{} accurately and concretely is crucial. Once given an \downtaskname{}, seeking a \methodname{} that directly optimizes the \taskname{} can provide \usernames{} with straightforward answers. 
In summary, the goals of interpretability are not impossible to achieve, but they may require new \methodnames{} along with a precise understanding of the assumptions under which such \methodnames{} are reliable. 
%
The development of \methodnames{} that come equipped with performance guarantees and enjoy computational efficiency remains a major open challenge for future work.


\clearpage